\section*{Chapter \ref{ch:g10:algebra} --- Algebra: Equations and Inequalities}

\begin{solution}{exo:g10:algebra:1}
$3(2x-5) - 2(x-7) = 6x - 15 - 2x + 14 = 4x - 1$.

$(x+4)(2x-1) = 2x^2 - x + 8x - 4 = 2x^2 + 7x - 4$.

$(3x-2)^2 = 9x^2 - 12x + 4$.

$(5-x)(5+x) = 25 - x^2$.
\end{solution}

\begin{solution}{exo:g10:algebra:2}
$x^2 - 49 = (x+7)(x-7)$ (difference of squares).

$x^2 + 10x + 25 = (x+5)^2$.

$7x^2 - 14x = 7x(x - 2)$ (common factor).

$16x^2 - 8x + 1 = (4x - 1)^2$.
\end{solution}

\begin{solution}{exo:g10:algebra:3}
$4x + 9 = x - 3$: $3x = -12$, so $x = -4$.

$\frac{2x-1}{3} = 5$: $2x - 1 = 15$, so $x = 8$.

$2(x-3) = 2x + 1$ expands to $2x - 6 = 2x + 1$, i.e.\ $-6 = 1$: false for
every $x$. No solution --- the two sides define parallel lines that never
meet.
\end{solution}

\begin{solution}{exo:g10:algebra:4}
$(x+5)(2x-8) = 0$: $x = -5$ or $x = 4$.

$x^2 - 16 = (x-4)(x+4) = 0$: $x = 4$ or $x = -4$.

$x^2 + 6x + 9 = (x+3)^2 = 0$: the single solution $x = -3$.
\end{solution}

\begin{solution}{exo:g10:algebra:5}
$3x - 5 \leq 7$: $3x \leq 12$, so $x \leq 4$: solution set
$\intoc{-\infty}{4}$.

$-4x + 1 > 9$: $-4x > 8$, dividing by $-4$ reverses the inequality:
$x < -2$: solution set $\intoo{-\infty}{-2}$.

$2(x-1) \geq 5x + 4$: $2x - 2 \geq 5x + 4$, so $-3x \geq 6$, so
$x \leq -2$: solution set $\intoc{-\infty}{-2}$.
\end{solution}

\begin{solution}{exo:g10:algebra:6}
The common factor is $(x-1)$:
\[
E = (x-1)\bigl[(x+3) - (2x - 5)\bigr] = (x-1)(x + 3 - 2x + 5)
= (x-1)(8 - x).
\]
Then $E = 0$ when $x = 1$ or $x = 8$ (zero-product rule).
\end{solution}

\begin{solution}{exo:g10:algebra:7}
$x^2 = 5x$: $x^2 - 5x = x(x - 5) = 0$, so $x = 0$ or $x = 5$ (never
divide by $x$: the solution $0$ would be lost).

$(2x+1)^2 = (x-4)^2$: bring to the left and factor the difference of
squares:
\[
(2x+1)^2 - (x-4)^2
= \bigl[(2x+1) + (x-4)\bigr]\bigl[(2x+1) - (x-4)\bigr]
= (3x - 3)(x + 5).
\]
Zero-product rule: $x = 1$ or $x = -5$.
\end{solution}

\begin{solution}{exo:g10:algebra:8}
$(x+2)(4-x) > 0$: the factor $x + 2$ vanishes at $-2$ (positive after),
$4 - x$ vanishes at $4$ (positive before). The product is positive
exactly when both factors are positive, i.e.\ on $\intoo{-2}{4}$.

$\frac{2x-6}{x+1} \geq 0$: numerator zero at $3$, denominator zero at
$-1$ (forbidden). \omterm{def:g10:algebra:signtable}{Sign table}: the quotient is positive when both parts
have the same sign, i.e.\ for $x < -1$ or $x > 3$; it vanishes at $x = 3$.
Solution set: $\intoo{-\infty}{-1} \cup \intco{3}{+\infty}$.
\end{solution}

\begin{solution}{exo:g10:algebra:9}
After $x$ months the first service costs $9x + 15$ and the second $12x$.
The first is cheaper when
\[
9x + 15 < 12x
\ \Longleftrightarrow\
15 < 3x
\ \Longleftrightarrow\
x > 5 .
\]
From the $6$th month on, the first service is the cheaper choice overall.
\end{solution}

\begin{solution}{exo:g10:algebra:10}
Forbidden value: $x = 1$. For $x \neq 1$, multiply both sides by $x - 1$:
\[
x + 3 = 2(x - 1) = 2x - 2
\ \Longleftrightarrow\
x = 5 .
\]
Since $5 \neq 1$, the solution is $x = 5$. Check:
$\frac{5+3}{5-1} = \frac84 = 2$.
\end{solution}

\begin{solution}{exo:g10:algebra:11}
$(n+1)^2 - n^2 = n^2 + 2n + 1 - n^2 = 2n + 1$, which is odd for every
\omterm{def:g10:numbers:sets}{integer} $n$.

Two consecutive odd numbers can be written $2n - 1$ and $2n + 1$. Then
\[
(2n+1)^2 - (2n-1)^2
= \bigl[(2n+1)+(2n-1)\bigr]\bigl[(2n+1)-(2n-1)\bigr]
= 4n \times 2 = 8n,
\]
a multiple of $8$.
\end{solution}

\begin{solution}{pb:g10:algebra:1}
\textbf{1.} $x^2 - 10x + 25 = (x - 5)^2$;
$9x^2 - 49 = (3x - 7)(3x + 7)$;
$x^2 + 6x + 9 - 4 = (x + 3)^2 - 2^2 = (x + 1)(x + 5)$.

\textbf{2.} $x = -1$ or $x = -5$; \quad $x = \frac73$ or
$x = -\frac73$.

\textbf{3.} $(5 - t)(5 + t) = 25 - t^2 = 21$ gives $t^2 = 4$,
$t = 2$: the numbers are $3$ and $7$ (sum $10$, product $21$:
check).

\textbf{4.} $(7 - t)(7 + t) = 49 - t^2 = 33$: $t = 4$, numbers
$3$ and $11$.

\textbf{5.} $\left(\frac s2 - t\right)\left(\frac s2 + t\right)
= \left(\frac s2\right)^2 - t^2 = p$ forces
$t^2 = \left(\frac s2\right)^2 - p$. A solution exists exactly
when this is $\geq 0$, i.e.\ when
$p \leq \left(\frac s2\right)^2$ --- the product of two numbers
of given sum can never beat the square of the half-sum
(queen Dido's fence problem, in the Middle School volume knew).

\textbf{6.} $(x - a)(x - b) = x^2 - (a + b)x + ab$: the
\omterm{def:g10:algebra:equation}{equation} $x^2 - sx + p = 0$ says precisely ``$x$ is one of two
numbers with sum $s$ and product $p$''.

\textbf{7.} $(x-5)^2 - 4 = x^2 - 10x + 25 - 4 = x^2 - 10x +
21$. Difference of squares:
$(x - 5 - 2)(x - 5 + 2) = (x - 7)(x - 3) = 0$: solutions $3$
and $7$ --- the scribe's numbers of question 3.

\textbf{8.} $x^2 + 6x - 7 = (x + 3)^2 - 16 = 0$ gives
$(x + 3)^2 = 16$, $x + 3 = \pm 4$: $x = 1$ or $x = -7$.

\textbf{9.} $x^2 + 4x + 5 = (x + 2)^2 + 1 \geq 1 > 0$ for every
$x$: no real solution. In the form $(x + h)^2 + k$: solutions
exist exactly when $k \leq 0$ (then $(x+h)^2 = -k$ has the two
roots $-h \pm \sqrt{-k}$); when $k > 0$ the expression stays
positive.

\textbf{10.} The square of side $x$ (area $x^2$) with two
$3 \times x$ rectangles ($6x$ in all) forms an L of area
$x^2 + 6x = 7$; the missing $3 \times 3$ corner (area $9$)
completes a big square of side $x + 3$. Adding $9$ to both
sides: $(x + 3)^2 = 16$, so $x + 3 = 4$ (lengths are positive)
and $x = 1$. The algebra and the picture are the same act.

\textbf{11.} $(x-3)(x+2)$: negative between the roots. Solution
of $< 0$: $x \in \intoo{-2}{3}$.

\textbf{12.} $x^2 - 10x + 21 \leq 0$, i.e.\
$(x - 3)(x - 7) \leq 0$: between the roots, ends included:
$x \in \intcc{3}{7}$.

\textbf{13.} Zero at $x = 1$, forbidden at $x = -2$; the
quotient is positive outside the roots:
$x \in \intoo{-\infty}{-2} \cup \intco{1}{+\infty}$.

\textbf{14.} $25 - (x - 5)^2 = 25 - x^2 + 10x - 25 =
x(10 - x)$: identity verified. Since $(x-5)^2 \geq 0$ always,
$A(x) \leq 25$, with equality exactly when $x = 5$: the
$5 \times 5$ square pen, maximum $25$ m$^2$ --- the grade-6
table, now a two-line proof.

\textbf{15.} For $x > 0$, multiply the claim by $x$:
$x^2 + 1 \geq 2x$, i.e.\ $x^2 - 2x + 1 = (x - 1)^2 \geq 0$ ---
true, with equality only at $x = 1$. Dividing back by $x > 0$
preserves everything (\cref{prop:g10:algebra:ineqrules}).

\textbf{16.} Cross-multiplying $\frac1x = \frac{x}{1 - x}$
(all lengths positive): $x^2 = 1 - x$, i.e.\
$x^2 + x - 1 = 0$.

\textbf{17.} $x^2 + x - 1 = \left(x + \frac12\right)^2 -
\frac54 = 0$ gives $x + \frac12 = \pm\frac{\sqrt5}{2}$. The
positive root: $x = \frac{\sqrt5 - 1}{2} \approx 0.618$.

\textbf{18.} $\varphi = \frac{2}{\sqrt5 - 1} =
\frac{2(\sqrt5 + 1)}{(\sqrt5 - 1)(\sqrt5 + 1)}
= \frac{2(\sqrt5 + 1)}{5 - 1} = \frac{\sqrt5 + 1}{2}
\approx 1.618$.

\textbf{19.} $\varphi^2 = \frac{(\sqrt5 + 1)^2}{4} =
\frac{6 + 2\sqrt5}{4} = \frac{3 + \sqrt5}{2}$, and
$\varphi + 1 = \frac{\sqrt5 + 3}{2}$: equal. And
$\varphi - 1 = \frac{\sqrt5 - 1}{2} = x = \frac1\varphi$ (that
is how $\varphi$ was defined). The golden ratio is the number
whose square adds one and whose reciprocal subtracts one.

\textbf{20.} $\sqrt5$ is \omterm{ex:g10:numbers:classify}{irrational} (the irrationality weekend problem of the Middle School volume: $5$
is no perfect square), and adding $1$ then halving preserves
irrationality (rational operations, \cref{pb:g10:numbers:1}).
Ratios: $\frac53 \approx 1.667$, $\frac85 = 1.600$,
$\frac{13}{8} = 1.625$, $\frac{21}{13} \approx 1.615$:
alternately above and below, closing in on
$\varphi = 1.618\ldots$ --- the rhythm-counting numbers of
the rhythm-counting weekend problem of the Middle School volume secretly chase the golden section; the
chapters on sequences will catch them in the act.
\end{solution}
